% Appendix F -- Theory (body). Source: main/reports/E4_theory.md (derivation 2, predictions 3, grid 4, results 5,
%   mismatches 6); numbers from main/results/E4/e4_comparisons.json and manifest.json (keys f.*, e4.*).
This appendix summarizes a derivation under a model: Gaussian sampling planners on a locally quadratic cost, with the quantization error treated as a random field. It was written, and its simulations were run on CPUs, before the E1 results were read.

\textbf{Setting.} In this appendix $J$ denotes a cost, to be minimized (Section~\ref{sec:setup} writes $J$ for a return). In one evaluation (planner iteration $t$) the planner draws candidates $a_1,\dots,a_N$ from its population, with true cost $J_i$, error $e_i$ and model cost $\hat J_i=J_i+e_i$. Its weights $w_i$ depend only on the model costs and are shift invariant, $W(\hat J+c\mathbf 1)=W(\hat J)$, as are argmin, the top-$K$ mean of CEM and the softmax over the top $K$ of MPPI; the next mean is $\sum_i w_ia_i$. The OEI of Section~\ref{sec:oei} is the identity $\mathrm{OEI}=\mathrm{sel}_{\mathrm{last}}+\mathrm{drift}$.

\textbf{Selection within one evaluation.} If $(J_i,e_i)$ are i.i.d.\ bivariate normal with standard deviations $\sigma_J,\sigma_e$ and correlation $\rho$, then for every such rule
\[
\mathbb E\Big[\textstyle\sum_i w_ie_i\Big]-\mathbb E[\bar e]=-\frac{\sigma_e^2+\rho\sigma_J\sigma_e}{\sigma_{\hat J}}\,S_{\mathrm{op}},\qquad S_{\mathrm{op}}=\mathbb E\Big[\textstyle\sum_i w_iZ_i\Big],\quad Z_i=\frac{\mu_{\hat J}-\hat J_i}{\sigma_{\hat J}},
\]
exactly, since $e_i$ splits into a multiple of $\hat J_i$ and a part independent of the weights. $S_{\mathrm{op}}$ is the selection intensity of the rule: about $\sqrt{2\ln N}$ for argmin, nearly constant in $N$ for a top-$K$ mean with a fixed elite fraction, and growing with $N$ when $K$ is fixed. Without the normality of $J$, if $e$ is normal with covariance $\Sigma_e$ and independent of $J$, Stein's identity gives the selection term as $\sum_{i,j}(\Sigma_e)_{ij}\,\mathbb E[\partial w_i/\partial\hat J_j]$; shift invariance makes every error component common to all candidates cancel exactly, and only the per-candidate variance $\sigma_\perp^2$ remains, multiplied by the density of candidates at the selection boundary. Selection therefore exploits error in proportion to the density of near ties, and common random numbers remove the shared part. Redrawn noise produces this winner's curse as well: scored with the draw that selected it, the executed plan's error is negative in expectation, and scored with an independent draw it is zero. This predicted, before E1, that TD-MPC2's OEI, which uses the same draw, would be negative under \shiter, as it was (Appendices~\ref{app:mech} and~\ref{app:neg}).

\textbf{Redrawn unbiased noise does not drift.} If the error of iteration $t$ is $m(a)+\varepsilon_t(a)$ with $\varepsilon_t$ independent of everything before iteration $t$, the last population's positions are independent of $\varepsilon_I$, so $\mathbb E[\mathrm{drift}]=\mathbb E[\bar m_{p_I}-\bar m_{p_1}]$, which is zero for first-order unbiased noise ($m$ constant), for any $N$, $K$, $I$ and planner; a draw shared by the candidates adds a common component that contributes variance, not bias. What remains is a second-order field: for unbiased dither the expected model cost differs from the unquantized one by $b_2(a)=\tfrac12\sum_\ell\Delta_\ell^2\sum_{i\in\ell}\mathcal H_{ii}(a)v_i(a)$, with $\mathcal H$ the downstream Hessian and $v_i=f_i(1-f_i)$ for stochastic rounding ($f_i$ the fractional part) or $\tfrac1{12}$ for \detshift. It is deterministic in expectation and hence exploitable, but it scales as $4^{-k}$ while first-order rounding error scales as $2^{-k}$.

\textbf{A persistent field is exploited.} With $e(a)=b(a)$ fixed across iterations, the elite update follows, to first order, the gradient of the model cost smoothed by the population,
\[
\mu_{t+1}-\mu_t\approx-\frac{S_{\mathrm{op}}\sigma_t^2}{\sigma_{\hat J,t}}\nabla\hat J_{\sigma_t}(\mu_t),\qquad \hat J_\sigma(\mu)=\mathbb E_z\big[\hat J(\mu+\sigma z)\big]=J_\sigma+b_\sigma,
\]
so the iterations converge to the minimizer of $J+b$ smoothed at the final population width: the planner solves the perturbed problem exactly, which is the bias of sample-average approximation~\citep{kleywegt2002saa}. For a weak stationary Gaussian field of variance $\sigma_b^2$ and correlation length $\lambda$ around a quadratic $J$ with Hessian $\mathcal H$ in $d$ dimensions,
\[
\mathbb E[\mathrm{drift}]\approx-\frac{\sigma_b^2\,\mathrm{tr}\,\mathcal H^{-1}}{\lambda^2}\Big(1+\frac{2\sigma_I^2}{\lambda^2}\Big)^{-d/2-1}=-2\,\mathbb E[\mathrm{regret}].
\]
The drift does not depend on $N$ to first order at a fixed elite fraction, grows with the iterations until it saturates, and is not averaged away by iterating. A field that is rough on the population's scale ($\lambda\ll\sigma_t$) is smoothed out and acts between iterations like redrawn noise, so the population schedule decides how much of a deterministic rounding field a planner exploits: MPPI's population width stops at its floor $\sigma_{\min}$, while CEM's keeps shrinking.

\textbf{Fixed dither and \detshift.} A dither table kept for a plan or an episode is a persistent field of the same kind, and its exploitable variance across candidates has an exact expression in the table's sharing pattern; it reduces to $A/N_{\mathrm{eff}}+B$ only when the units sharing an offset move together, so the table-size law is not guaranteed in a real network (Appendix~\ref{app:mech}). The error of \detshift is the sawtooth part of the stochastic-rounding error (subtractive dither): uniform with variance $\Delta^2/12$ for every input and redrawn at every call. It belongs to the redrawn-shared class, has no drift and less exploitable variance than \shiter, so the theory predicts $D_{\detshift}\le D_{\shiter}$ at equal bits.

\textbf{Simulation.} The two planners were re-implemented from their code (TD-MPC2's MPPI with 512 candidates, 64 elites, 6 iterations, a standard-deviation range of $[0.05, 2]$, 24 policy proposals and a sampled elite; DINO-WM's CEM with 300 candidates, top 30 and 30 iterations) and run on quadratic costs of 2 to 12 dimensions, with seven error models: a persistent random-feature field, redrawn i.i.d.\ noise, a redrawn shared field with a common shock, and toy rounding models under \dtm, \shiter, a fixed table and \detshift. A pre-declared grid of \nv{f.jobs} jobs varied $N$, the elite rule, the error scale, the table size, the bit width and the number of iterations, with 128 problem instances and two planner seeds per cell, for \nv{e4.nruns} planner runs in all. Table~\ref{tab:theory} compares the derivation with the simulation.

\begin{table}[h]
\caption{Derivation against simulation (E4). Slopes are from regressions of measured on predicted values through the origin; ``cells'' are grid cells, each the mean of 256 planner runs.}
\label{tab:theory}
\centering\small
\setlength{\tabcolsep}{4pt}
\begin{tabular}{@{}>{\raggedright}p{4.3cm}>{\raggedright}p{3.6cm}>{\raggedright\arraybackslash}p{6.4cm}@{}}
\toprule
Quantity & Prediction & Simulation \\
\midrule
Selection, one evaluation, normal $(J,e)$ & exact formula & median relative error \nv{f.c1.err}\% over \nv{f.c1.n} cells \\
Selection term inside the planners, same draw & formula with the rule's $S_{\mathrm{op}}$ & slope \nv{f.c2.elite.slope} ($R^2$ \nv{f.c2.elite.r2}; CEM elite mean, \nv{f.c2.elite.n} cells), \nv{f.c2.argmin.slope} (\nv{f.c2.argmin.r2}; CEM argmin, near-tie count), \nv{f.c2.mppi.slope} (\nv{f.c2.mppi.r2}; MPPI, \nv{f.c2.mppi.n} cells) \\
Redrawn noise: drift and re-scored selection & zero & intervals contain 0 in \nv{f.c3.fresh.cover}\% (drift) and \nv{f.c3.redraw.cover}\% (selection) of cells, nominal 95\%; same-draw selection below 0 in \nv{f.c3.same.below}\% \\
Persistent field: drift & negative; depth from the smoothed problem & negative in \nv{f.c3.pers.cem}\% (CEM) and \nv{f.c3.pers.mppi}\% (MPPI) of cells; depth slope \nv{f.c4.slope}, $R^2$ \nv{f.c4.r2} \\
Fixed table: variance across candidates & exact formula & measured/exact \nv{f.c9.ratio}, $R^2$ \nv{f.c9.r2} \\
\detshift vs.\ \shiter, equal bits & $D_{\detshift}\le D_{\shiter}$ & significantly worse in \nv{f.c10.worse} of \nv{f.c10.of} cells; median damage ratio \nv{f.c10.med} \\
Second-order dither field $b_2$ & $\propto4^{-k}$ ($-2$ in $\log_2$ per bit) & \nv{f.c11.slope} per bit \\
\bottomrule
\end{tabular}
\end{table}

\textbf{What the theory did not get right.} (i) A persistence index, the agreement of plans across re-runs, does not distinguish persistent from redrawn error (medians \nv{f.c12.persistent} against \nv{f.c12.fresh}), because a strong persistent field has several local minima and different seeds reach different ones; drift is the better measure, and the experiments use it. (ii) The smoothed problem predicts how deep the planner exploits a persistent field but not which minimum it reaches. (iii) The steady-state damage of redrawn independent noise under MPPI has the predicted shape but not its size or its dependence on the number of elites. (iv) In the simulated CEM, which has no standard-deviation floor, iterating did not average out redrawn noise, because the collapsing population freezes the displacement it causes; in DINO-WM, however, \shiter stayed at the unquantized level for every number of CEM iterations tested (E2). (v) The predicted growth of the benefit of sharing with the near-tie density did not appear in the experiments (H4b, Appendix~\ref{app:neg}), and the $1/N_{\mathrm{eff}}$ law held only partly (Appendix~\ref{app:mech}). The predictions that held are the negative OEI of redrawn noise in TD-MPC2's definition, the flat drift of a persistent field in $N$, the growth of $|\mathrm{OEI}|$ with $N$ only at a fixed elite count (H2c), \detshift being no worse than \shiter at equal bits, and, under MPPI, damage of redrawn noise that falls with the iterations and equals that of a persistent error at one iteration (Section~\ref{sec:persist}).
